% Fichier engendré par tools/generate-tex.py à partir de 03-suites.
% Ne pas éditer à la main : tout le texte provient des commentaires du fichier Lean.
\documentclass[11pt,a4paper]{article}

% Compile aussi bien avec pdflatex qu'avec xelatex ou tectonic.
\usepackage{iftex}
\ifPDFTeX
  \usepackage[T1]{fontenc}
  \usepackage[utf8]{inputenc}
\else
  \usepackage{fontspec}
\fi
\usepackage[french]{babel}
\usepackage{amsmath,amssymb,amsthm}
\usepackage[margin=2.5cm]{geometry}
\usepackage[hidelinks]{hyperref}

% Les figures sont dessinées en TikZ : elles reprennent ainsi les
% polices du document. Voir la skill `illustrate-theorem`.
\usepackage{tikz}
\usetikzlibrary{calc}

\theoremstyle{definition}
\newtheorem{definition}{Définition}
\theoremstyle{plain}
\newtheorem{theoreme}{Théorème}
\newtheorem{lemme}[theoreme]{Lemme}

% Un énoncé écrit mais non démontré : la preuve Lean est un `sorry`.
\newcommand{\admis}{\par\smallskip\noindent\textit{Admis : l'énoncé est écrit en Lean, sa démonstration reste à faire.}\par}

% Renvoi à la déclaration Lean, une ligne après la démonstration, aligné à droite.
\newcommand{\source}[2]{\par\smallskip\noindent\hfill{\footnotesize\href{#1}{\texttt{#2}}}\par\smallskip}

\title{Suites}
\date{}

\begin{document}
\maketitle
% Les identifiants Lean en machine à écrire ne se coupent pas : on tolère des
% espacements plus lâches plutôt que des lignes qui débordent.
\sloppy

\section{Suites}

\noindent
Lycée \textemdash{} section \og{} Suites \fg{}.
Une suite est une application de \(\mathbb{N}\) dans \(\mathbb{R}\). Les limites sont celles de Mathlib :
\texttt{Tendsto u atTop (nhds \(\ell\))} pour la convergence vers \texttt{\(\ell\)}, \texttt{Tendsto u atTop atTop} pour la
divergence vers \texttt{+∞}. Les deux se déplient en la définition avec \texttt{\(\varepsilon\)} et avec \texttt{A} du
programme, ce que le premier énoncé vérifie.
Énoncés et démonstrations en français : voir Suites.tex.

\subsection{Définitions}

\noindent
Toutes les définitions du fichier sont posées ici, avant les démonstrations.

\begin{definition}
La \emph{suite arithmétique} de premier terme $u_0$ et de raison $r$ est
$u_n = u_0 + nr$.
\end{definition}
\source{https://github.com/Commutator-IO/learn-lean/blob/main/courses/02-lycee/03-suites/Suites.lean\#L20}{Suites.lean\#L20}

\begin{definition}
La \emph{suite géométrique} de premier terme $u_0$ et de raison $q$ est
$u_n = u_0 q^n$.
\end{definition}
\source{https://github.com/Commutator-IO/learn-lean/blob/main/courses/02-lycee/03-suites/Suites.lean\#L23}{Suites.lean\#L23}

\begin{definition}
Deux suites sont \emph{adjacentes} lorsque l'une croît, l'autre décroît, et que leur
différence tend vers zéro.
\end{definition}
\source{https://github.com/Commutator-IO/learn-lean/blob/main/courses/02-lycee/03-suites/Suites.lean\#L27}{Suites.lean\#L27}

\subsection{Suites arithmétiques et géométriques : terme général}

\begin{theoreme}
Terme général d'une suite arithmétique : $u_n = u_0 + nr$.
\end{theoreme}

\begin{proof}
C'est la définition, écrite sous forme explicite.
\end{proof}
\source{https://github.com/Commutator-IO/learn-lean/blob/main/courses/02-lycee/03-suites/Suites.lean\#L33}{Suites.lean\#L33}

\begin{theoreme}
Une suite arithmétique passe d'un terme au suivant en ajoutant la raison :
$u_{n+1} = u_n + r$.
\end{theoreme}

\begin{proof}
On compare les deux termes généraux :
\[ u_0 + (n+1)r = \bigl(u_0 + nr\bigr) + r . \]
\end{proof}
\source{https://github.com/Commutator-IO/learn-lean/blob/main/courses/02-lycee/03-suites/Suites.lean\#L37}{Suites.lean\#L37}

\begin{theoreme}
Terme général d'une suite géométrique : $u_n = u_0 q^n$.
\end{theoreme}

\begin{proof}
C'est la définition, écrite sous forme explicite.
\end{proof}
\source{https://github.com/Commutator-IO/learn-lean/blob/main/courses/02-lycee/03-suites/Suites.lean\#L44}{Suites.lean\#L44}

\begin{theoreme}
Une suite géométrique passe d'un terme au suivant en multipliant par la raison :
$u_{n+1} = u_n \times q$.
\end{theoreme}

\begin{proof}
On compare les deux termes généraux :
\[ u_0 q^{n+1} = \bigl(u_0 q^n\bigr) \times q . \]
\end{proof}
\source{https://github.com/Commutator-IO/learn-lean/blob/main/courses/02-lycee/03-suites/Suites.lean\#L48}{Suites.lean\#L48}

\subsection{Sens de variation}

\begin{theoreme}
Une suite arithmétique de raison positive est croissante.
\end{theoreme}

\begin{proof}
Si $m \le n$ alors $mr \le nr$ puisque $r \ge 0$, donc $u_0 + mr \le u_0 + nr$. C'est le
signe de la raison qui décide du sens de variation.
\end{proof}
\source{https://github.com/Commutator-IO/learn-lean/blob/main/courses/02-lycee/03-suites/Suites.lean\#L56}{Suites.lean\#L56}

\begin{theoreme}
De raison négative, elle est décroissante.
\end{theoreme}

\begin{proof}
Même calcul : multiplier l'inégalité $m \le n$ par $r \le 0$ la renverse.
\end{proof}
\source{https://github.com/Commutator-IO/learn-lean/blob/main/courses/02-lycee/03-suites/Suites.lean\#L64}{Suites.lean\#L64}

\begin{theoreme}
Une suite géométrique de premier terme positif et de raison au moins un est croissante :
c'est le signe de $q - 1$ qui décide.
\end{theoreme}

\begin{proof}
La suite des puissances $q^n$ est croissante dès que $q \ge 1$ ; la multiplier par
$u_0 \ge 0$ conserve l'ordre.
\end{proof}
\source{https://github.com/Commutator-IO/learn-lean/blob/main/courses/02-lycee/03-suites/Suites.lean\#L73}{Suites.lean\#L73}

\begin{theoreme}
De raison comprise entre zéro et un, elle est décroissante.
\end{theoreme}

\begin{proof}
Les puissances d'un nombre de $[0, 1]$ décroissent, et la multiplication par
$u_0 \ge 0$ conserve l'ordre.
\end{proof}
\source{https://github.com/Commutator-IO/learn-lean/blob/main/courses/02-lycee/03-suites/Suites.lean\#L81}{Suites.lean\#L81}

\subsection{Limite d'une suite : la définition avec \texttt{\(\varepsilon\)}}

\begin{theoreme}
La convergence de la bibliothèque est exactement la définition du programme :
\[ u_n \to \ell \iff \forall \varepsilon > 0,\ \exists N,\ \forall n \ge N,\
   |u_n - \ell| < \varepsilon . \]
\end{theoreme}

\begin{proof}
La convergence, dans la bibliothèque, n'est pas définie avec $\varepsilon$ mais avec des
\emph{filtres} : $u$ tend vers $\ell$ lorsque l'image du filtre « rangs assez grands » est
plus fine que le filtre des voisinages de $\ell$. Autrement dit : tout voisinage de $\ell$
contient les termes à partir d'un certain rang.

L'équivalence avec la définition scolaire tient à ce que, dans un espace métrique, les
boules forment une base de voisinages : tout voisinage de $\ell$ contient une boule
$B(\ell, \varepsilon)$, et toute boule est un voisinage. Se donner un voisinage revient donc
à se donner un $\varepsilon > 0$, et
\[ u_n \in B(\ell, \varepsilon) \iff d(u_n, \ell) < \varepsilon
   \iff |u_n - \ell| < \varepsilon , \]
la dernière écriture étant la distance de $\mathbb{R}$.

La démonstration ne fait que cela : citer la caractérisation métrique, puis déplier la
distance. Les deux définitions ne sont pas seulement équivalentes, elles sont la même,
énoncée à deux niveaux de généralité.
\end{proof}
\source{https://github.com/Commutator-IO/learn-lean/blob/main/courses/02-lycee/03-suites/Suites.lean\#L91}{Suites.lean\#L91}

\begin{figure}[h]
\centering
% Construction : axe des rangs de 0 à 8, limite ℓ = 1.6, bande ℓ ± 0.45.
%   Termes u_n aux rangs 1..8, alternant autour de ℓ en s'en rapprochant :
%   (1, 0.4) (2, 2.6) (3, 0.95) (4, 2.25) (5, 1.3) (6, 1.9) (7, 1.5) (8, 1.72).
%   Rang N = 5 : au-delà, tous les termes sont dans la bande.
% SVG : x = 30 + 30x, y = 130 - 40y.
\begin{tikzpicture}[scale=1, >=stealth, every node/.style={font=\small}]
  \draw[->] (0, 0) -- (9, 0) node[right] {$n$};
  \draw[dashed] (0, 1.6) -- (8.8, 1.6) node[right] {$\ell$};
  \draw[dotted] (0, 2.05) -- (8.8, 2.05) node[right] {$\ell + \varepsilon$};
  \draw[dotted] (0, 1.15) -- (8.8, 1.15) node[right] {$\ell - \varepsilon$};
  \foreach \n/\v in {1/0.4, 2/2.6, 3/0.95, 4/2.25, 5/1.3, 6/1.9, 7/1.5, 8/1.72}
    \fill (\n, \v) circle (1.6pt);
  \draw[|-] (5, -0.35) -- (5, -0.05);
  \node[below] at (5, -0.35) {$N$};
\end{tikzpicture}

\caption{La définition de la limite. Une fois la bande d'amplitude $\varepsilon$ choisie autour de $\ell$, il existe un rang $N$ au-delà duquel tous les termes y sont \textemdash{} et cela pour tout $\varepsilon$, si petit soit-il.}
\end{figure}

\begin{theoreme}
Et la divergence vers $+\infty$ est la définition avec $A$ :
\[ u_n \to +\infty \iff \forall A,\ \exists N,\ \forall n \ge N,\ A \le u_n . \]
\end{theoreme}

\begin{proof}
Même vérification, les voisinages de $+\infty$ étant les intervalles $[A ; +\infty[$.
\end{proof}
\source{https://github.com/Commutator-IO/learn-lean/blob/main/courses/02-lycee/03-suites/Suites.lean\#L99}{Suites.lean\#L99}

\subsection{Unicité de la limite}

\begin{theoreme}
Unicité de la limite : une suite a au plus une limite.
\end{theoreme}

\begin{proof}
Si la suite tendait vers deux limites distinctes, les voisinages disjoints de celles-ci
devraient tous deux contenir les termes à partir d'un certain rang, ce qui est impossible.
Le programme admet ce résultat ; il tient à la séparation de la droite réelle.
\end{proof}
\source{https://github.com/Commutator-IO/learn-lean/blob/main/courses/02-lycee/03-suites/Suites.lean\#L106}{Suites.lean\#L106}

\subsection{Opérations sur les limites}

\begin{theoreme}
La limite d'une somme est la somme des limites.
\end{theoreme}

\begin{proof}
Si $|u_n - \ell| < \varepsilon/2$ et $|v_n - \ell'| < \varepsilon/2$ au-delà d'un
certain rang, alors $|(u_n + v_n) - (\ell + \ell')| < \varepsilon$ par inégalité
triangulaire.
\end{proof}
\source{https://github.com/Commutator-IO/learn-lean/blob/main/courses/02-lycee/03-suites/Suites.lean\#L113}{Suites.lean\#L113}

\begin{theoreme}
Celle d'un produit est le produit des limites.
\end{theoreme}

\begin{proof}
Même principe, en majorant $|u_n v_n - \ell\ell'|$ par
$|u_n||v_n - \ell'| + |\ell'||u_n - \ell|$, le premier facteur étant borné puisque la
suite converge.
\end{proof}
\source{https://github.com/Commutator-IO/learn-lean/blob/main/courses/02-lycee/03-suites/Suites.lean\#L118}{Suites.lean\#L118}

\begin{theoreme}
Celle d'un quotient est le quotient des limites, à condition que le dénominateur ne
tende pas vers zéro : c'est là que naissent les formes indéterminées.
\end{theoreme}

\begin{proof}
L'hypothèse $\ell' \ne 0$ garantit que le dénominateur ne s'annule pas au voisinage de
l'infini, ce qui permet de majorer l'écart des quotients.
\end{proof}
\source{https://github.com/Commutator-IO/learn-lean/blob/main/courses/02-lycee/03-suites/Suites.lean\#L124}{Suites.lean\#L124}

\begin{theoreme}
Forme indéterminée : deux suites tendant toutes deux vers $+\infty$ peuvent avoir une
différence de limite quelconque — ici $c$, choisi à l'avance.
\end{theoreme}

\begin{proof}
Les suites $n + c$ et $n$ tendent toutes deux vers $+\infty$, et leur différence est
constante égale à $c$. « $\infty - \infty$ » n'a donc pas de valeur : c'est ce qu'on
appelle une forme indéterminée.
\end{proof}
\source{https://github.com/Commutator-IO/learn-lean/blob/main/courses/02-lycee/03-suites/Suites.lean\#L131}{Suites.lean\#L131}

\subsection{Limite de \texttt{qⁿ}}

\begin{theoreme}
Si $|q| < 1$, la suite $q^n$ tend vers zéro.
\end{theoreme}

\begin{proof}
Les puissances d'un nombre de module strictement inférieur à un décroissent
géométriquement vers zéro.
\end{proof}
\source{https://github.com/Commutator-IO/learn-lean/blob/main/courses/02-lycee/03-suites/Suites.lean\#L142}{Suites.lean\#L142}

\begin{theoreme}
Si $q > 1$, elle tend vers $+\infty$.
\end{theoreme}

\begin{proof}
L'inégalité de Bernoulli donne $q^n \ge 1 + n(q-1)$, qui tend vers $+\infty$.
\end{proof}
\source{https://github.com/Commutator-IO/learn-lean/blob/main/courses/02-lycee/03-suites/Suites.lean\#L147}{Suites.lean\#L147}

\begin{theoreme}
Si $q = 1$, elle est constante égale à un.
\end{theoreme}

\begin{proof}
$1^n = 1$ pour tout $n$.
\end{proof}
\source{https://github.com/Commutator-IO/learn-lean/blob/main/courses/02-lycee/03-suites/Suites.lean\#L152}{Suites.lean\#L152}

\subsection{Théorèmes de comparaison et théorème des gendarmes}

\begin{theoreme}
Théorème des gendarmes : une suite encadrée par deux suites de même limite converge
vers cette limite.
\end{theoreme}

\begin{proof}
Si $v_n \le u_n \le w_n$ et si $v_n$ et $w_n$ tendent toutes deux vers $\ell$, alors
au-delà d'un certain rang les deux encadrants sont dans l'intervalle
$]\ell - \varepsilon ; \ell + \varepsilon[$, donc $u_n$ aussi.
\end{proof}
\source{https://github.com/Commutator-IO/learn-lean/blob/main/courses/02-lycee/03-suites/Suites.lean\#L160}{Suites.lean\#L160}

\begin{figure}[h]
\centering
% Construction : axe des rangs de 0 à 8, limite commune ℓ = 1.5.
%   v_n monte vers ℓ  : 0.5, 0.85, 1.1, 1.25, 1.34, 1.4, 1.44, 1.47
%   w_n descend vers ℓ : 2.5, 2.15, 1.9, 1.75, 1.66, 1.6, 1.56, 1.53
%   u_n reste entre les deux, sans monotonie :
%                        1.15, 1.95, 1.2, 1.7, 1.4, 1.55, 1.47, 1.51
%   Chaque famille est reliée par une ligne brisée, et nommée à gauche, là où
%   les trois sont encore écartées.
% SVG : x = 30 + 32x, y = 130 - 40y.
\begin{tikzpicture}[scale=1, >=stealth, every node/.style={font=\small}]
  \draw[->] (0, 0) -- (9, 0) node[right] {$n$};
  \draw[dashed] (0.4, 1.5) -- (8.7, 1.5) node[right] {$\ell$};
  \draw[gray!60] plot coordinates
    {(1,0.5) (2,0.85) (3,1.1) (4,1.25) (5,1.34) (6,1.4) (7,1.44) (8,1.47)};
  \draw[gray!60] plot coordinates
    {(1,2.5) (2,2.15) (3,1.9) (4,1.75) (5,1.66) (6,1.6) (7,1.56) (8,1.53)};
  \draw[gray!60] plot coordinates
    {(1,1.15) (2,1.95) (3,1.2) (4,1.7) (5,1.4) (6,1.55) (7,1.47) (8,1.51)};
  \foreach \n/\v in {1/0.5, 2/0.85, 3/1.1, 4/1.25, 5/1.34, 6/1.4, 7/1.44, 8/1.47}
    \fill (\n, \v) circle (1.5pt);
  \foreach \n/\v in {1/2.5, 2/2.15, 3/1.9, 4/1.75, 5/1.66, 6/1.6, 7/1.56, 8/1.53}
    \fill (\n, \v) circle (1.5pt);
  \foreach \n/\v in {1/1.15, 2/1.95, 3/1.2, 4/1.7, 5/1.4, 6/1.55, 7/1.47, 8/1.51}
    \draw[fill=white] (\n, \v) circle (2pt);
  \node[left] at (0.85, 2.5) {$w_n$};
  \node[left] at (0.85, 1.15) {$u_n$};
  \node[left] at (0.85, 0.5) {$v_n$};
\end{tikzpicture}

\caption{Le théorème des gendarmes. La suite $u_n$, prise entre $v_n$ et $w_n$ qui convergent toutes deux vers $\ell$, n'a plus la place d'aller ailleurs.}
\end{figure}

\begin{theoreme}
Comparaison : une suite minorée par une suite qui tend vers $+\infty$ tend vers
$+\infty$.
\end{theoreme}

\begin{proof}
Si $v_n \le u_n$ et si $v_n$ dépasse tout $A$ à partir d'un rang, alors $u_n$ le dépasse
aussi.
\end{proof}
\source{https://github.com/Commutator-IO/learn-lean/blob/main/courses/02-lycee/03-suites/Suites.lean\#L167}{Suites.lean\#L167}

\subsection{Convergence monotone}

\begin{theoreme}
Convergence monotone : toute suite croissante majorée converge.
\end{theoreme}

\begin{proof}
L'ensemble des valeurs de la suite est non vide et majoré : il admet donc une borne
supérieure, vers laquelle la suite converge, la croissance assurant qu'elle s'en approche
autant qu'on veut.

Le programme admet ce théorème. Il n'est pas anodin : c'est une forme de la propriété de
la borne supérieure, c'est-à-dire de la complétude de $\mathbb{R}$ — ce qui distingue les
réels des rationnels.
\end{proof}
\source{https://github.com/Commutator-IO/learn-lean/blob/main/courses/02-lycee/03-suites/Suites.lean\#L175}{Suites.lean\#L175}

\begin{theoreme}
Une suite croissante non majorée tend vers $+\infty$.
\end{theoreme}

\begin{proof}
Soit $A$ un réel. La suite n'étant pas majorée, l'un de ses termes $u_n$ dépasse $A$ ;
la croissance fait que tous les suivants aussi. C'est la définition de la divergence vers
$+\infty$.
\end{proof}
\source{https://github.com/Commutator-IO/learn-lean/blob/main/courses/02-lycee/03-suites/Suites.lean\#L180}{Suites.lean\#L180}

\subsection{Toute suite convergente est bornée}

\begin{theoreme}
Toute suite convergente est bornée.
\end{theoreme}

\begin{proof}
Au-delà d'un certain rang, tous les termes sont à distance moins de un de la limite,
donc bornés ; avant ce rang, il n'y a qu'un nombre fini de termes, dont on prend le plus
grand en valeur absolue.
\end{proof}
\source{https://github.com/Commutator-IO/learn-lean/blob/main/courses/02-lycee/03-suites/Suites.lean\#L190}{Suites.lean\#L190}

\subsection{Suites adjacentes}

\begin{theoreme}
Deux suites adjacentes convergent vers une même limite.
\end{theoreme}

\begin{proof}
Montrons d'abord que $u_n \le v_n$ pour tout $n$. Sinon, pour un certain rang $n$ on
aurait $v_n < u_n$, et pour tout $m \ge n$ la croissance de $u$ et la décroissance de $v$
donneraient $v_m - u_m \le v_n - u_n < 0$ ; en passant à la limite, on obtiendrait
$0 \le v_n - u_n$, contradiction.

La suite $u$ est donc croissante et majorée par $v_0$ : elle converge vers un réel $\ell$.
Comme $v_n = (v_n - u_n) + u_n$ et que la différence tend vers zéro, $v$ converge vers le
même $\ell$.

La dichotomie, second volet de l'énoncé du programme, n'est pas traitée ici : elle demande
de construire la suite des intervalles emboités et d'en déduire l'existence d'une racine,
ce qui relève du théorème des valeurs intermédiaires.
\end{proof}
\source{https://github.com/Commutator-IO/learn-lean/blob/main/courses/02-lycee/03-suites/Suites.lean\#L198}{Suites.lean\#L198}

\subsection{Suites définies par récurrence \texttt{uₙ₊₁ = f(uₙ)}}

\begin{theoreme}
Suites définies par récurrence : si $u_{n+1} = f(u_n)$, si $f$ est continue et si
$u_n \to \ell$, alors $f(\ell) = \ell$.
\end{theoreme}

\begin{proof}
La suite décalée $u_{n+1}$ tend vers la même limite $\ell$ que $u_n$. Par continuité de
$f$, la suite $f(u_n)$ tend vers $f(\ell)$. Or ces deux suites sont la même, et une suite
n'a qu'une limite : $f(\ell) = \ell$.

C'est ce qui permet de trouver la limite en résolvant l'équation $f(x) = x$ — à condition
d'avoir établi \emph{par ailleurs} que la suite converge, ce que ce théorème ne dit pas.
\end{proof}
\source{https://github.com/Commutator-IO/learn-lean/blob/main/courses/02-lycee/03-suites/Suites.lean\#L223}{Suites.lean\#L223}

\subsection{Énoncés admis}

\noindent
Ce que le programme demande et que ce document ne démontre pas. Les énoncés sont écrits en Lean, pour que le manque se compte ; leur démonstration est admise.

\begin{theoreme}
Deux suites adjacentes \textemdash{} l'une croissante, l'autre décroissante, leur
différence tendant vers zéro \textemdash{} convergent vers une même limite.
\end{theoreme}
\admis
\source{https://github.com/Commutator-IO/learn-lean/blob/main/courses/02-lycee/03-suites/Suites.lean\#L239}{Suites.lean\#L239}

\vfill
\noindent\rule{0.35\linewidth}{0.4pt}\par
\noindent{\footnotesize Leonardo de Moura et Sebastian Ullrich,
\emph{The Lean~4 Theorem Prover and Programming Language},
\emph{Automated Deduction — CADE~28}, Springer, 2021, p.~625--635,
\href{https://doi.org/10.1007/978-3-030-79876-5_37}{doi:10.1007/978-3-030-79876-5\_37}.\par}

\end{document}
